\documentclass[11pt,a4paper]{article}

% --- Eingabe / Sprache ---
\usepackage{fontspec}                 % XeTeX/tectonic: natives UTF-8, korrektes ß/ä/ö/ü
\usepackage[ngerman]{babel}
\addto\captionsngerman{\renewcommand{\abstractname}{Abstract (English)}}
\usepackage{microtype}

% --- Anführungszeichen ---
% Zitatkonvention dieses Papers: Im deutschen Fliesstext durchgaengig \glqq...\grqq
% (globale Stilvorgabe). \enquote{} wird AUSSCHLIESSLICH im englischen Abstract
% verwendet (bewusste EN-Ausnahme via csquotes/babel). Kein Mischbetrieb im Body.
\usepackage[autostyle=true,german=quotes]{csquotes}

% --- Mathematik ---
\usepackage{amsmath}
\usepackage{amssymb}
\usepackage{amsthm}

% --- Tabellen ---
\usepackage{booktabs}
\usepackage{array}
\usepackage{longtable}

% --- Code-Listings (F#/Lean) ---
\usepackage{xcolor}
\usepackage{listings}

% --- Verweise (cleveref muss nach hyperref geladen werden) ---
\usepackage[hidelinks]{hyperref}
\usepackage[capitalise,ngerman]{cleveref}

% --- Einheitliche Auszeichnung der A/B-Verifizierbarkeitsmarken ---
\newcommand{\AB}[1]{\textbf{#1}}

% --- Theorem-Umgebungen ---
\theoremstyle{definition}
\newtheorem{definition}{Definition}[section]
\newtheorem{satz}{Satz}[section]
\newtheorem{lemma}{Lemma}[section]
\newtheorem{korollar}{Korollar}[section]
\theoremstyle{remark}
\newtheorem{bemerkung}{Bemerkung}[section]

% --- cleveref-Namen ---
\crefname{bemerkung}{Bemerkung}{Bemerkungen}
\Crefname{bemerkung}{Bemerkung}{Bemerkungen}
\crefname{satz}{Satz}{Saetze}
\Crefname{satz}{Satz}{Saetze}
\crefname{definition}{Definition}{Definitionen}
\Crefname{definition}{Definition}{Definitionen}
\crefname{lstlisting}{Listing}{Listings}
\Crefname{lstlisting}{Listing}{Listings}

% --- Umbruch langer Bezeichner/URLs (reduziert Overfull-\hbox) ---
\usepackage{xurl}
\emergencystretch=3em
\sloppy

% --- Listing-Stile für F# und Lean ---
\lstdefinelanguage{fsharp}{
  morekeywords={let,rec,match,with,type,module,open,member,override,
    namespace,and,if,then,else,for,in,do,fun,function,of,mutable,
    true,false,when,yield,return},
  sensitive=true,
  morecomment=[l]{//},
  morecomment=[s]{(*}{*)},
  morestring=[b]",
}
\lstdefinelanguage{lean}{
  morekeywords={theorem,lemma,def,by,intro,intros,simp,rw,exact,apply,
    have,show,fun,match,with,structure,inductive,namespace,open,end,
    omega,calc},
  sensitive=true,
  morecomment=[l]{--},
  morecomment=[s]{/-}{-/},
  morestring=[b]",
}
\lstset{
  basicstyle=\ttfamily\small,
  keywordstyle=\color{blue!60!black}\bfseries,
  commentstyle=\color{green!40!black}\itshape,
  stringstyle=\color{red!50!black},
  numbers=left,
  numberstyle=\tiny\color{gray},
  frame=single,
  breaklines=true,
  showstringspaces=false,
  captionpos=b,
  columns=fullflexible,
}

\title{Das Terminierungs-Orakel: Konvergenz agentischer
Softwareentwicklung gegen eine externalisierte, getypte Spezifikation}
\author{Cong Chanh Vinzenz Nguyen}
\date{19.\ Juni 2026}

\begin{document}
\maketitle

\begin{abstract}
\noindent
Agentic software development is increasingly organized as a loop: a generating
agent proposes, a checking step accepts or rejects, and the result feeds the next
iteration. We argue that the load-bearing problem of this loop is not the
generator and not the iteration, but its \emph{termination oracle} --- the
instance deciding when a proposal is \enquote{done}. A loop without a
forgery-resistant termination oracle does not converge towards correctness but
towards whatever its checking step happens to wave through. Our answer is not
better generation but \emph{convergence against an externalized, typed
specification} that carries the termination oracle. We instantiate this in a
method (\texttt{cong-driven-development}, a typed F\# discriminated-union graph
\texttt{SPOT} with the convergence states \texttt{Aligned}/\texttt{Pending}/%
\texttt{Diverged}/\texttt{Orphaned}) and two case studies: \texttt{runenruf}
(existence proof, $46/46$ tests, $\sim\!125$ LoC, the honestly friendliest domain
where the spec \emph{is} the world) and \texttt{ledger-casestudy} (counterpart
with real IO, a mutable requirement, $6/6$ tests, and two Lean-proved,
\texttt{sorry}-free invariants --- value preservation and non-negativity). We separate the
generator/oracle independence into two axes: a \emph{mechanical} one (tool
allowlist without generator write access, real and machine-checkable) and an
\emph{epistemic} one (decorrelation of spec author and code author, required only
when the spec stems from a separate source). Every load-bearing claim is declared
either \textbf{A} (machine- and externally verifiable: commit hash, arXiv id, CI
run, \texttt{sorry} audit) or \textbf{B} (a human setting). We make no priority
claim and prove no impossibility about machines: a second, sufficiently
independent agent could set the ground invariant. We close with a gap, not a
promise --- two small case studies, one proved invariant, no unbounded autonomy:
one screw turned further.
\end{abstract}

\tableofcontents
\bigskip

\section{Einleitung}
\label{sec:einleitung}

Agentische Softwareentwicklung organisiert sich zunehmend als Schleife: ein
generierender Agent erzeugt einen Vorschlag, ein Prüfschritt akzeptiert oder verwirft ihn,
und das Ergebnis fließt in die nächste Iteration ein. Boris Cherny bringt diese
Verschiebung auf den Punkt, wenn er die Tätigkeit des Entwicklers als \glqq my job is to
write loops\grqq{} beschreibt~\cite{cherny2026}. Wir argumentieren, dass das eigentliche
Problem dieser Schleife nicht der Generator und nicht die Iteration selbst ist, sondern ihr
\emph{Terminierungs-Orakel}: jene Instanz, die entscheidet, wann ein Vorschlag \glqq fertig\grqq{}
ist. Eine Schleife ohne fälschungssicheres Terminierungs-Orakel konvergiert nicht gegen
Korrektheit, sondern gegen das, was ihr Prüfschritt zufällig durchwinkt. Unser Lösungsansatz
ist entsprechend nicht eine bessere Generierung, sondern die \emph{Konvergenz gegen eine
externalisierte, getypte Spezifikation}, die das Terminierungs-Orakel trägt.

\paragraph{Cognitive Debt und Intent Debt.}
Die Notwendigkeit einer externalisierten Spezifikation ist nicht bloß ergonomisch. Storey
beschreibt einen Schuldbegriff jenseits der klassischen technischen Schuld:
\emph{Intent Debt} als \glqq the absence of externalized rationale that developers and AI
agents need to work safely with code\grqq{}~\cite{storey2026}. Fehlt die externalisierte
Absicht, so verlagert sich das Wissen darüber, \emph{warum} ein Code korrekt sein soll, in
nicht inspizierbare Zustände --- in den menschlichen Kopf oder in den verworfenen Kontext
eines Agenten. Lahiri formuliert die korrespondierende Aufgabe als benannte offene
\emph{Grand Challenge}: die Formalisierung von Intent als Voraussetzung verlässlicher
KI-gestützter Programmierung~\cite{lahiri2026}. Wir behandeln diese Externalisierung als
Konstruktionsprinzip, nicht als Nebenprodukt.

\paragraph{Ehrlicher Stand der Technik.}
Drei Befunde rahmen die Ausgangslage und zugleich ihre Grenzen. Erstens zur Reichweite
heutiger Agenten: METR misst einen \glqq 50\%-task-completion time horizon\grqq{} und berichtet
für Anfang 2025 eine Aufgabenlänge von etwa $50$ bis $110$ Minuten bei einer Verdopplung
von rund $7$ Monaten~\cite{metr2025}; wir übernehmen ausschließlich diese Zahlen und leiten
daraus keine Extrapolation ab. Zweitens zur prinzipiellen Schranke: nach dem Satz von
Rice ist jede nicht-triviale semantische Eigenschaft von Programmen unentscheidbar
\cite{rice1953}. Die Voraussetzung des Satzes --- eine nicht-triviale semantische
Eigenschaft --- ist für \glqq erfüllt die Spezifikation\grqq{} gegeben; daraus folgt aber nur eine
unentscheidbare \emph{obere Schranke} über alle Programme. Sie motiviert die entscheidbare,
\emph{pro Spezifikation} approximierende Prüfung, die wir verwenden, und ist insofern
\emph{Motivation}, nicht Existenzbeweis unserer Konstruktion. Wir betonen ausdrücklich: aus dem
Satz von Rice folgt \emph{nicht}, dass ein Mensch im Prüfvorgang unersetzlich wäre. Drittens
zur Vertrauenswürdigkeit von Spezifikationen selbst: Misu, Ma und Lopes zeigen mit
\textsc{VeriAct}, dass \emph{verifizierer-akzeptiert} nicht \emph{korrekt} bedeutet --- eine vom
Verifizierer angenommene formale Spezifikation kann die Absicht verfehlen~\cite{veriact2026}.
Dieser Befund ist die nächste Nachbarschaft unserer Arbeit und zugleich die Quelle ihres
ehrlichsten Vorbehalts.

\paragraph{Leitfrage.}
Daraus ergibt sich die Frage, die dieses Papier strukturiert: \emph{Lässt sich die
agentische Schleife so konstruieren, dass ihr Terminierungs-Orakel an einer
externalisierten, getypten Spezifikation hängt --- und welcher Rest bleibt nach maximaler
Mechanisierung notwendig menschlich?} Wir beantworten den ersten Teil konstruktiv anhand
einer Methode und zweier Fallstudien; den zweiten Teil beantworten wir, indem wir den
verbleibenden Rest exakt lokalisieren, statt ihn wegzuargumentieren.

\paragraph{A/B-Methodik dieses Papers.}
Wir unterwerfen jede tragende Behauptung einer expliziten Disziplin: sie deklariert sich
entweder als \AB{A} --- maschinell und extern verifizierbar (Commit-Hash, \texttt{arXiv}-ID,
CI-Lauf, \texttt{sorry}-Audit) --- oder als \AB{B} --- eine menschliche Setzung
(normative Wahl, Spezifikation-gegen-Welt-Bindung, Neuheit als Literaturaussage,
Generalisierung über Domänen). Keine Aussage bleibt unmarkiert. Diese Markierung ist selbst
der zentrale Hebel der Arbeit: Sie trennt, was Logik leistet, von dem, was der Mensch
einbringt, und macht den \glqq Rest\grqq{} prüfbar statt rhetorisch. \Cref{sec:methode} legt die
Methode offen; die A/B-Tabelle in \Cref{sec:grenzen-ab} bindet jede Marke an ihren Anker.
Wir kündigen damit zugleich an, dass dieses Papier \emph{mit} einer Lücke argumentiert: zwei
kleine Fallstudien, zwei Lean-bewiesene Invarianten, keine unbeschränkte Autonomie ---
eine Schraube weiter, nicht das gelöste Problem.

\section{Methode}
\label{sec:methode}

\subsection{Konvergenz gegen eine externalisierte Spezifikation}
\label{sec:konvergenz}

Den Kern bildet eine Umkehrung des Akzeptanzkriteriums. Statt \glqq der Agent hält an, wenn er
zufrieden ist\grqq{} gilt \glqq der Agent hält an, wenn ein externes, getyptes Orakel die
Konvergenz gegen die Spezifikation bestätigt\grqq. Die Methode und ihre IDE-Realisierung liegen
in \texttt{cong-driven-development}~\cite{cdd2026}, ausgewertet am Commit \texttt{main}. Das tragende Artefakt
ist \texttt{SPOT} (\emph{Single Point of Truth}): ein typisierter F\#-Discriminated-Union-Graph,
in dem jeder Knoten als ein \texttt{JSON} pro Knoten unter \texttt{.spot/} materialisiert ist.
Der Konvergenzzustand jedes Knotens ist auf vier Werte getypt: \texttt{Aligned},
\texttt{Pending}, \texttt{Diverged}, \texttt{Orphaned}. Das Orakel
\texttt{SetzeSpecAligned} setzt einen Test-Knoten genau dann auf \texttt{Aligned}, wenn seine
Id als Marker im Test\-quellcode vorkommt (\texttt{covered}); die Funktion prüft Marker-Präsenz,
keinen Testlauf. Dass diese Marker zu bestandenen Tests gehören, erzwingt die grüne CI-Suite
(rot blockt den Merge) zusammen mit dem reflexiven Selbst-Test --- daraus folgt \texttt{Aligned}
$\Rightarrow$ Test-PASS, nicht bloßer Exit-Code $0$. Das Selbst-Modell der
Methode umfasst $66$ Knoten, davon $62$ \texttt{Aligned} und $4$ \texttt{Pending}, bei
$37/37$ grünen Tests (autoritativer Runner-Count @\,\texttt{main}); der reflexive
Selbst-Gate trägt den Knoten \texttt{spec-gate-selbst-hart}. Dass das Gate auf sich selbst
angewandt grün ist, ist eine A-Aussage; dass es das \emph{Richtige} prüft, ist eine B-Aussage
(siehe \Cref{bem:spec-welt}).

\subsection{Generator-/Orakel-Trennung --- zweiachsig}
\label{sec:trennung}

Die Trennung von Erzeuger und Prüfer ist nicht eine einzige Eigenschaft, sondern zerfällt in
zwei voneinander unabhängige Achsen.

\paragraph{Mechanische Achse (Generator $\perp$ Prüf-Orakel).}
Der Generator besitzt kein Schreibrecht auf das Prüf-Orakel. Technisch erzwungen wird dies
durch eine Werkzeug-Allowlist ohne Schreibzugriff des Generators auf den Gate-Pfad; der
Generator kann das Urteil über seine eigene Ausgabe nicht fälschen. Diese Achse ist real und
maschinell prüfbar --- sie ist eine A-Aussage, verankert im Code @\,\texttt{main}.

\paragraph{Epistemische Achse (Generator $\perp$ Spezifikation-Autor).}
Die zweite Achse betrifft die \emph{Herkunft} der Spezifikation. Sie gilt nur, sofern die
Spezifikation aus einer von der Codegenerierung \emph{separaten} Quelle stammt. Erzeugt
dieselbe Modell-Instanz sowohl Spezifikation als auch Code, so korrelieren ihre Fehler ---
ein \emph{common-mode}-Versagen, genau der von \textsc{VeriAct} beschriebene Fall, in dem die
verifizierer-akzeptierte Spezifikation mit dem Code gemeinsam die Absicht
verfehlt~\cite{veriact2026}. Auf dieser Achse ist der Mensch der \emph{Dekorrelator}: nicht
weil eine Maschine die Rolle prinzipiell nicht einnehmen könnte --- ein zweiter, unabhängiger
Agent könnte die Spezifikation ebenso stellen ---, sondern weil eine unabhängige Quelle
\emph{erforderlich} ist und in unseren Fallstudien der Mensch diese Quelle war. Konkret wurde
der in \texttt{ledger-casestudy} gefangene Fehlschluss --- ein Modell, das sich selbst als
\glqq Falsifiable, after 4 tests\grqq{} auswies --- dadurch gefangen, dass ein Mensch eine stärkere
Property und ein geändertes Requirement als Gegen-Orakel einbrachte; nicht die Schleife aus
sich selbst heraus. Die Notwendigkeit der Dekorrelation ist strukturell (A: die Korrelation ist
beobachtbar); \emph{wer} dekorreliert, ist eine Setzung (B).

\begin{bemerkung}[Spezifikation gegen Welt]
\label{bem:spec-welt}
Die mechanische Achse garantiert, dass der Code die Spezifikation erfüllt; sie garantiert
nicht, dass die Spezifikation die \emph{Welt} trifft. Die Bindung \glqq Spezifikation $\models$
Welt\grqq{} ist nach \textsc{VeriAct} nicht orakel-prüfbar~\cite{veriact2026} und damit
durchgängig eine B-Aussage. Wir behaupten an keiner Stelle Spezifikation-$\leftrightarrow$-Code-Äquivalenz.
\end{bemerkung}

\subsection{Verschiebung des Akzeptanzkriteriums}
\label{sec:akzeptanz}

Die Methode verschiebt das Akzeptanzkriterium von einer impliziten, generator-internen
Zufriedenheit zu einem expliziten, externen, getypten Prädikat. Diese Verschiebung ist die
operationale Antwort auf das Terminierungs-Orakel-Problem aus \Cref{sec:einleitung}: Die
Schleife terminiert genau dann, wenn der \texttt{SPOT}-Graph keinen \texttt{Diverged}- und
keinen \texttt{Orphaned}-Knoten mehr trägt und jeder relevante Knoten über
\texttt{SetzeSpecAligned} auf \texttt{Aligned} gehoben ist. Der Loop selbst bleibt unverändert
agentisch; verändert ist allein, woran sein Haltepunkt hängt. Im Sinne von
\textsc{LLM-Modulo}~\cite{kambhampati2024} ist der LLM der \emph{Generator}, das getypte
Orakel der \emph{Kritiker}; im Sinne property-orientierten Feedbacks~\cite{pgs2025} liefert der
\texttt{Diverged}-Zustand eine strukturell minimale, property-bezogene Rückmeldung statt eines
diffusen \glqq versuche es erneut\grqq.

\subsection{Abgrenzung}
\label{sec:abgrenzung}

Spezifikationsgetriebene Werkzeuge wie Kiro oder GitHub Spec-Kit externalisieren Absicht
ebenfalls, jedoch typischerweise als \emph{natürlichsprachliche} Spezifikationsdokumente, deren
Erfüllung nicht über ein getyptes, ausführbares Orakel mit fälschungssicherer
Generator-/Prüfer-Trennung erzwungen wird. Der Unterschied liegt nicht im Vorhandensein einer
Spezifikation, sondern in deren Typisierung, ihrer maschinellen Prüfbarkeit und der
architektonisch erzwungenen Trennung nach \Cref{sec:trennung}. Wir erheben hieraus
\emph{keinen} Prioritätsanspruch: Uns ist keine publizierte Implementierung mit genau dieser
Kombination bekannt --- getypte externalisierte Spezifikation, architektonisch getrennter
Generator und Orakel, Werkzeug-Allowlist ohne Generator-Schreibrecht sowie ein
property- und beweis-verifiziertes Gate ---, doch diese Neuheit ist eine \emph{Literaturaussage}
(B), kein Beweis. Sie kann durch einen einzigen Gegenbeleg fallen.

\begin{bemerkung}[Setzen der Grund-Invariante]
\label{bem:grundinvariante}
Eine Konsequenz der Konstruktion verdient explizite Trennung in eine strukturelle und eine
normative Aussage. \emph{(i) Strukturell (Regress):} Die Grund-Invariante, gegen die geprüft
wird, ist im System nicht selbst-fundierbar --- jede Begründung einer Invariante setzt eine
weitere Invariante voraus; der Regress terminiert nur durch einen Anker \emph{außerhalb} des
Prüfvorgangs. Dies ist ein architektonisches Argument, kein im Papier geführter formaler
Beweis. \emph{(ii) Normativ:} Dass der \emph{Mensch} diesen Anker setzt, ist eine
\emph{Wahl} --- begründet darin, dass die Absicht extern zum Generator und
verantwortungsbehaftet ist ---, und \emph{keine} Unmöglichkeitsaussage über Maschinen. Ein
zweiter Agent könnte die Invariante setzen; der Mensch ist hier \emph{verschoben}, nicht durch
ein Theorem als notwendig erwiesen. Wir vermeiden bewusst jede Formulierung im Sinne von
\glqq unmechanisierbar\grqq{} oder \glqq nur ein Mensch kann\grqq{}; das wäre der Penrose/Lucas-Fehlschluss,
den wir ablehnen. Die Theoreme aus \Cref{sec:einleitung} --- Rice als Motivation, Gödel und
Tarski (sofern überhaupt herangezogen) als bloße Analogie --- tragen diese Aussage nicht: aus
keinem von ihnen folgt die Unersetzlichkeit eines Menschen.
\end{bemerkung}

\section{Fallstudie A: Runenruf}
\label{sec:runenruf}

Runenruf ist der \emph{Existenzbeweis} der Methode: die Frage, ob eine getypte,
externalisierte Spezifikation als Terminierungs-Orakel eines agentischen Loops überhaupt
einen sorry-freien, deterministisch reproduzierbaren Konvergenz-Zustand erzeugen kann, wird
hier an einer minimalen Domäne positiv beantwortet. Wir markieren die Domänenwahl sofort als
das, was sie ist: Runenruf ist die \emph{ehrlich freundlichste} Domäne, die wir kennen --- eine
reine Simulationsdomäne, in der die Spezifikation mit der Welt zusammenfällt
(\Cref{bem:runenruf-grenze}). Das Artefakt belegt Anwendbarkeit, nicht Skalierung.

\subsection{Aufbau}
\label{sec:runenruf-aufbau}

Der Code liegt in \texttt{github.com/Koschnag/runenruf}, ausgewertet am Commit
\texttt{56376a5}. Die Simulationsdomäne umfasst rund $125$ Zeilen
F\#/.NET-Code\,@\,\texttt{56376a5}. Über dem Generator-Loop steht das in
\Cref{sec:konvergenz} beschriebene Orakel \texttt{SetzeSpecAligned},
das einen Test-Knoten auf \texttt{Aligned} setzt, sobald ein Test-Marker ihn im
Quellcode abdeckt; die CI-erzwungene grüne Suite bindet dieses \texttt{Aligned} an einen
bestandenen Test --- gewertet wird das Test-PASS, nicht der Prozess-Exitcode $0$.

\subsection{Verifikation}
\label{sec:runenruf-verifikation}

Der autoritative Test-Runner meldet $46/46$ grüne Tests\,@\,\texttt{56376a5}; die Zahl stammt
aus der Runner-Ausgabe, nicht aus einem \texttt{grep} über Testnamen. Tragend ist eine
FsCheck-Property mit dem Spec-Schlüssel \texttt{spec-siegel-lager-nichtnegativ}: sie
quantifiziert über \emph{jeden} Seed und \emph{jede} Folge von Domänenoperationen und behauptet
die Nichtnegativität des Siegel-Lagerbestands. Die Auswertung ist deterministisch --- gleicher
Seed erzeugt gleichen Lauf ---, sodass ein Gegenbeispiel reproduzierbar bleibt und nicht von
einem nichtdeterministischen Scheduler abhängt.

\begin{bemerkung}[Verifizierbarkeitsachse]
\label{bem:runenruf-achse}
Die Aussage \glqq $46/46$ grün\grqq\ ist eine \AB{A}-Behauptung im Sinne von
\Cref{sec:einleitung}: maschinell extern verifiziert, an \texttt{github.com/Koschnag/runenruf}\,@\,\texttt{56376a5}
und den dortigen CI-Run gebunden. Die FsCheck-Property ist ebenfalls \AB{A}: sie wird vom
Runner ausgeführt. Was \emph{nicht} aus diesen Läufen folgt, ist die Bindung der Spezifikation
an eine externe Welt --- diese fällt in \AB{B}.
\end{bemerkung}

\begin{bemerkung}[Grenze: die Spec ist die Welt]
\label{bem:runenruf-grenze}
In Runenruf existiert kein externer Referent jenseits der Simulation. Die Spezifikation
\emph{ist} die Welt, daher entfällt die Spec-gegen-Welt-Bindung
(\glqq{}Spec\,$\models$\,Welt\grqq), die nach VeriAct~\cite{veriact2026} nicht
orakel-prüfbar ist und in unserer Klassifikation zu \AB{B} gehört. Genau deshalb ist
Runenruf ein Existenzbeweis und kein Realismus-Beweis: er zeigt, dass die
Typen-Property-Maschinerie geschlossen \emph{durchläuft}, nicht, dass sie etwas außerhalb ihrer
selbst trifft. Der härtere Fall mit echtem IO und veränderlichem Requirement folgt in
\Cref{sec:ledger}.
\end{bemerkung}

\section{Fallstudie B: Ledger}
\label{sec:ledger}

Ledger ist das \emph{Gegenstück} zu Runenruf: dieselbe Methode in einer Domäne mit echtem IO,
einem über die Zeit \emph{veränderlichen} Requirement und einer geschlossenen Drei-Schichten-Kette
Typen\,$\rightarrow$\,FsCheck\,$\rightarrow$\,Lean-Beweis (\Cref{sec:ledger-schichten}). Hier
fällt die Spec nicht mit der Welt zusammen, und genau hier zeigt sich, wo das Loop-interne Orakel
endet und der menschliche Dekorrelator beginnt (\Cref{sec:ledger-misspec}).

\subsection{Aufbau}
\label{sec:ledger-aufbau}

Der Code liegt in \texttt{github.com/Koschnag/ledger-casestudy}, ausgewertet am Commit
\texttt{0f39ab0}; F\#/.NET~9, Beträge als \texttt{int64}-Cent, um Gleitkomma-Fehler an der
Wertgrenze auszuschließen. Anders als in \Cref{sec:runenruf} gibt es echtes IO: ein persistentes
Journal und einen Replay-Pfad. Die Kerndomäne umfasst rund $41$ Zeilen
F\#-Code\,@\,\texttt{0f39ab0}. Der autoritative Runner meldet $6/6$ grüne
Tests\,@\,\texttt{0f39ab0}.

\subsection{Gefangener Mis-Spec --- der Mensch als Gegen-Orakel}
\label{sec:ledger-misspec}

Im \texttt{.spot/}-Ledger findet sich der Eintrag \glqq Falsifiable, after 4 tests\grqq. Wir
attribuieren diesen Fund präzise: Er wurde \emph{nicht} aus dem Loop heraus erzeugt. Die ersten,
schwächeren Properties akzeptierten eine fehlerhafte Buchungslogik; verifizierer-akzeptiert ist
eben nicht korrekt~\cite{veriact2026}. Gefangen wurde der Fehler erst, als ein \emph{Mensch}
eine stärkere Property \emph{und} das geänderte Gebühren-Requirement als Gegen-Orakel einbrachte
--- erst diese externe Setzung machte die bis dahin grüne Spezifikation falsifizierbar. Das ist
der zweiachsige Befund aus \Cref{sec:trennung}: Generator\,$\perp$\,Prüf-Orakel ist mechanisch
erzwungen, aber Generator\,$\perp$\,Spec-Autor gilt nur bei separater Quelle. Erzeugt dieselbe
Modell-Instanz Spezifikation \emph{und} Code, korrelieren die Fehler (common-mode); der Mensch
ist hier der Dekorrelator, nicht der Loop aus sich selbst.

Der Befund ist als reproduzierbares Negativ-Artefakt konserviert:
\texttt{demo-gate-at-failure.sh}\,@\,\texttt{0f39ab0} führt das Gate in den Falsifikationszustand
und macht den Mis-Spec wiederholbar sichtbar. Anschließend wurde die Domäne auf doppelte
Buchführung (Doppik) umgestellt, sodass jede Buchung wertneutral bleibt.

\begin{bemerkung}[Verifizierbarkeitsachse]
\label{bem:ledger-achse}
\glqq $6/6$ grün\grqq, der sorry-freie Lean-Lauf (\Cref{sec:ledger-schichten}) und die Existenz
des Negativ-Artefakts sind \AB{A}-Behauptungen, gebunden an
\texttt{github.com/Koschnag/ledger-casestudy}\,@\,\texttt{0f39ab0}. Die Aussage, \emph{wer} die
stärkere Property setzt, und dass die Doppik-Spec die buchhalterische Welt trifft, sind
\AB{B}: normativ bzw.\ Spec-gegen-Welt, nicht aus dem Lauf ableitbar.
\end{bemerkung}

\subsection{Drei Schichten: Typen, FsCheck, Lean}
\label{sec:ledger-schichten}

Ledger trägt die einzige in dieser Arbeit \emph{geschlossene} Drei-Schichten-Kette:

\begin{enumerate}
  \item \textbf{Typen.} \texttt{int64}-Cent und die F\#-Datentypen schließen ganze
  Fehlerklassen (Gleitkomma, negative Repräsentationslücken) bereits statisch aus.
  \item \textbf{FsCheck.} Properties quantifizieren über Buchungsfolgen und prüfen
  Nichtnegativität, Isolation und Replay-Treue über zufällig generierte Eingaben.
  \item \textbf{Lean-Beweis.} \texttt{proofs/Werterhaltung.lean}\,@\,\texttt{0f39ab0} beweist die
  Werterhaltung für \emph{jede} Buchungsfolge --- ein Allquantor, der über das stichprobenhafte
  FsCheck hinausgeht.
\end{enumerate}

Entscheidend für die Aussagekraft des Lean-Beweises ist, dass das Lean-Modell den
Deckungs-/Ablehnungspfad \emph{enthält} und damit modellgleich zur F\#-Funktion \texttt{buche}
ist: eine Buchung, die die Deckung verletzen würde, wird im Modell abgelehnt, genau wie im Code.
Ein Beweis über ein Modell ohne Ablehnungspfad wäre vakuum-stark und würde die kritische
Eigenschaft am Code vorbeibeweisen.

\begin{lstlisting}[language=lean,caption={Kern-Theorem in \texttt{proofs/Werterhaltung.lean}\,@\,\texttt{0f39ab0} (verbatim; Definitionen und das Lemma \texttt{werterhaltung\_einzel} im Repo)},label={lst:werterhaltung}]
theorem werterhaltung (hb : Hauptbuch) (bs : List Buchung) :
    total (buchen hb bs) = total hb := by
  unfold buchen
  induction bs generalizing hb with
  | nil => rfl
  | cons b bs ih =>
      simp only [List.foldl_cons]
      rw [ih (buche hb b), werterhaltung_einzel]
\end{lstlisting}

\noindent
Der Beweis (\Cref{lst:werterhaltung}) ist \emph{sorry-frei}: \texttt{\#print axioms werterhaltung} liefert ausschließlich
\texttt{propext} und \texttt{Quot.sound}, also keine zusätzlichen Axiome jenseits des
Lean-Kerns. Der Beweis kommt ohne Mathlib aus (nur \texttt{Core}~$+$~\texttt{omega}); der
CI-Job \texttt{lean-proof} ist grün\,@\,\texttt{0f39ab0}.

\begin{bemerkung}[Was Lean hier \emph{nicht} abdeckt]
\label{bem:ledger-nur-werterhaltung}
Zwei Invarianten --- Werterhaltung und Nichtnegativität --- sind Lean-bewiesen.
Isolation und Replay-Treue sind ausschließlich durch FsCheck abgesichert. Die
Drei-Schichten-Kette Typen\,$\rightarrow$\,Property\,$\rightarrow$\,Beweis ist derzeit für
Werterhaltung und Nichtnegativität geschlossen; für die übrigen Invarianten fehlt die Beweis-Schicht. Wir rahmen
das ausdrücklich als Lücke, eine Schraube weiter, nicht als vollständige Verifikation der
Domäne.
\end{bemerkung}

\subsection{Kompression}
\label{sec:ledger-kompression}

Über beide Fallstudien hinweg ist die verifizierte Invariante deutlich kürzer als ihre
Implementierung: das Verhältnis liegt zwischen rund $40\times$ und $125\times$
(\Cref{sec:runenruf}: $\sim\!125$ LoC Domäne; \Cref{sec:ledger}: $\sim\!41$ LoC Domäne, jeweils
gegen die kompakte Property bzw.\ das Kern-Theorem). Wir formulieren die Lesart dieser Zahlen
streng konditional und eng.

\begin{bemerkung}[Kompression ist Diskrimination, keine Äquivalenz]
\label{bem:kompression}
Die Invariante \emph{fixiert} bzw.\ diskriminiert die kritische Korrektheitseigenschaft ---
Werterhaltung, Nichtnegativität ---, sie \emph{spezifiziert die Funktion nicht vollständig}.
Unendlich viele Implementierungen erfüllen dieselbe Invariante; aus
\glqq Code\,$\models$\,Invariante\grqq\ folgt keine Spec-Code-Äquivalenz. Die genannten
Faktoren gelten konditional auf genau diese zwei Fallstudien (N\,$=$\,$2$) und sind kein Beleg
für eine domänenübergreifende Kompressionsrate. Die Größenangaben sind \AB{A}
(Code\,@\,\texttt{56376a5} bzw.\ \texttt{0f39ab0}); der Schluss, dass die fixierte Eigenschaft
\glqq die richtige\grqq\ ist, ist \AB{B} (Spec-gegen-Welt).
\end{bemerkung}

\section{Reflexive Selbstanwendung}
\label{sec:reflexiv}

Eine Methode, die Korrektheit gegen eine externalisierte, getypte Spezifikation
erzwingt, muss sich der Frage stellen, ob sie auf sich selbst anwendbar ist. Wir
prüfen das nicht rhetorisch, sondern operativ: Die Methode \texttt{CDD} ist in
ihrem eigenen \texttt{.spot/}-Graphen modelliert, und ihr Gate läuft über sich
selbst.

\subsection{Selbst-Gate}
\label{sec:reflexiv-selbstgate}

Das Selbstmodell von \texttt{CDD}
ist im Repository \texttt{cong-driven-development}\,@\,\texttt{main} als typisierter
F\#-DU-Graph abgelegt: ein JSON-Knoten je Element unter \texttt{.spot/}, mit den
Konvergenzzuständen \texttt{Aligned}, \texttt{Pending}, \texttt{Diverged} und
\texttt{Orphaned}. Das Orakel \texttt{SetzeSpecAligned} markiert einen Test-Knoten
als \texttt{Aligned}, sobald ein Test-Marker ihn im Quellcode abdeckt; die grüne
CI-Suite (rot blockt den Merge) und der reflexive Selbst-Test binden dieses
\texttt{Aligned} an einen bestandenen Test, sodass \texttt{Aligned} den Test-\texttt{PASS}
impliziert, nicht den Exit-Code~0 des Prozesses. Diese Unterscheidung ist nicht kosmetisch:
Ein Prozess kann mit Exit-0 terminieren, ohne dass eine einzige Property tatsächlich geprüft wurde.

Der reflexive Selbst-Gate \texttt{spec-gate-selbst-hart}
(Repository \texttt{@main}) wendet exakt dieses Orakel auf das Selbstmodell
an. Der autoritative Testlauf (Runner, nicht \texttt{grep}) meldet
$37/37$~Tests grün; das Selbstmodell umfasst $66$~Knoten, davon $62$~\texttt{Aligned}
und $4$~\texttt{Pending}. Wir behaupten hier keine Vollständigkeit: Die vier
\texttt{Pending}-Knoten sind genau die Stellen, an denen die Methode ihre eigene
Spezifikation noch nicht eingeholt hat. Das ist der ehrliche, nicht der
geschönte Konvergenzzustand.

\subsection{Ehrliche Konvergenz}
\label{sec:reflexiv-konvergenz}

Konvergenz im Sinne von \texttt{CDD} ist kein Fixpunkt-Versprechen, sondern ein
beobachtbarer Zustand pro Knoten. Ein Graph mit $62$~\texttt{Aligned} und
$4$~\texttt{Pending} ist nicht \glqq fertig\grqq, sondern \emph{vermessen}: Die
Methode macht ihre eigene Lücke sichtbar, statt sie zu verbergen. Wir lesen das
als Funktionsmerkmal, nicht als Mangel --- ein Terminierungs-Orakel, das nur
\texttt{Aligned} kennt, wäre kein Orakel, sondern eine Zusicherung ohne
Diskriminationskraft.

\subsection{Die reflexive Grenze}
\label{sec:reflexiv-grenze}

Wendet man das Gate auf das Gate an, stellt sich die Frage nach der
\emph{Grund-Invariante}: Welche Invariante rechtfertigt das Orakel, das alle
übrigen Invarianten prüft? Hier sind zwei Aussagen strikt zu trennen.

\paragraph{(i) Strukturell --- der Regress.}
Die Grund-Invariante ist im System nicht selbst-fundierbar. Jede Begründung der
prüfenden Invariante $I_0$ müsste eine weitere Invariante $I_1$ heranziehen, die
$I_0$ als korrekt ausweist; deren Begründung wiederum $I_2$, und so fort. Der
Regress terminiert nicht durch ein internes Argument, sondern nur durch einen
Anker \emph{außerhalb} des Prüfvorgangs. Das ist eine strukturelle, architektonische
Aussage über das Prüfen --- ein Argument, kein im Papier geführter formaler Beweis,
und kein Werturteil.

\paragraph{(ii) Normativ --- die Wahl des Ankers.}
Dass im hier beschriebenen Aufbau der \emph{Mensch} diesen Anker setzt, ist eine
Wahl, keine Unmöglichkeitsaussage über Maschinen. Sie ruht auf zwei Gründen, die
beide normativ und nicht berechnungstheoretisch sind: Intent ist extern zum
Generator, und das Setzen des Ankers ist verantwortungsbehaftet. Ein zweiter,
hinreichend unabhängiger Agent könnte die Grund-Invariante ebenso setzen; der
Mensch ist dann \emph{verschoben}, nicht durch ein Theorem als notwendig
erwiesen. Wir behaupten ausdrücklich nicht, der Anker sei
\glqq unmechanisierbar\grqq{} oder \glqq nur ein Mensch\grqq{} könne ihn setzen ---
das wäre der Penrose/Lucas-Fehlschluss, den wir ablehnen
(vgl.\ \Cref{sec:reflexiv-theorie}).

\subsection{Theorie-Disclaimer}
\label{sec:reflexiv-theorie}
Drei metamathematische Resultate werden in dieser Arbeit in unterschiedlichen
Rollen genutzt, und die Rollen dürfen nicht vermischt werden. Der Satz von Rice
\cite{rice1953} ist ein \emph{Beweis}: Für jede nicht-triviale semantische
Eigenschaft partieller rekursiver Funktionen ist die zugehörige
Entscheidungsfrage unentscheidbar; die Voraussetzung (nicht-triviale semantische
Eigenschaft) ist im hier betrachteten Fall erfüllt. Rice liefert damit
ausschließlich die unentscheidbare obere Schranke und motiviert die
entscheidbare \emph{pro-Spec-Approximation} --- die Verschiebung von
\glqq korrekt schlechthin\grqq{} auf \glqq erfüllt diese eine getypte Spec\grqq.
Gödels Unvollständigkeitssätze und Tarskis Undefinierbarkeit dienen demgegenüber
nur als \emph{Motivation} und Analogie. Ihre Voraussetzungen --- rekursive
Axiomatisierung, Subsumtion der Robinson-Arithmetik~$Q$, Konsistenz,
Syntax-Kodierung --- beanspruchen wir für den \texttt{.spot/}-Graphen
\emph{nicht}, und sie gelten dort nicht. Aus keinem dieser drei Resultate folgt,
dass ein Mensch unersetzlich ist; der Regress (i) ist architektonisch, die
Anker-Setzung (ii) ist normativ.

\section{Grenzen und Geltungsbereich}
\label{sec:grenzen}

Wir halten die Grenzen dieser Arbeit explizit, weil ein Terminierungs-Orakel,
das seine eigene Reichweite überschätzt, denselben Fehler begeht, den es
verhindern soll.

\subsection{Empirische Basis}
\label{sec:grenzen-empirie}

Die Evidenz beruht auf \emph{zwei} kleinen Fallstudien.
\texttt{runenruf}\,@\,\texttt{56376a5} ($46/46$~Tests, Runner) ist mit $\sim 125$~LoC die
ehrlich freundlichste Domäne, weil die Spezifikation dort mit der Welt
zusammenfällt --- die Simulation \emph{ist} ihre eigene Realität.
\texttt{ledger-casestudy}\,@\,\texttt{0f39ab0} ($6/6$~Tests, Runner; F\#/.NET~9, int64-Cent,
$\sim 41$~LoC Domänenkern) ist das Gegenstück mit echtem IO (Journal und
Replay) und veränderlichem Requirement. Zwei Fallstudien zeigen
\emph{Anwendbarkeit}; sie beweisen keine Skalierung über Domänen ($N=2$). Es
gibt keinen Issue-Tracker im großen Stil und keine unbeschränkte Autonomie des
Loops.

\subsection{Reichweite der Verifikation}
\label{sec:grenzen-lean}

Im Repository \texttt{ledger-casestudy}\,@\,\texttt{0f39ab0} sind \emph{zwei} Invarianten
maschinell bewiesen: \texttt{proofs/Werterhaltung.lean} zeigt Werterhaltung und Nichtnegativität
für jede Buchungsfolge. Das Modell enthält den Deckungs- und Ablehnungspfad und
ist modellgleich zur F\#-Funktion \texttt{buche}; der Beweis ist
\texttt{sorry}-frei (\texttt{\#print axioms} weist nur \texttt{propext} und
\texttt{Quot.sound} aus), kommt ohne Mathlib aus (nur \texttt{Core} und
\texttt{omega}), und der CI-Job \texttt{lean-proof} ist grün
(\texttt{@0f39ab0}). Die übrigen Eigenschaften --- Isolation und
Replay-Treue --- sind \emph{nur} per FsCheck geprüft, etwa
\texttt{spec-siegel-lager-nichtnegativ} in \texttt{runenruf}\,@\,\texttt{56376a5} über jeden
Seed und jede Folge, deterministisch. Damit ist die Drei-Schichten-Kette
(Typen~$\to$~Property~$\to$~Beweis) derzeit für
Werterhaltung und Nichtnegativität geschlossen. Wir behaupten keine durchgehende Beweisabdeckung.

\subsection{Common-Mode-Fehler}
\label{sec:grenzen-commonmode}

Die mechanisch erzwungene Trennung Generator~$\perp$~Prüf-Orakel
(Werkzeug-Allowlist ohne Schreibrecht des Generators, fälschungssicher;
Code~\texttt{@main}) schützt nicht gegen jede Fehlerklasse. Wenn dieselbe
Modell-Instanz sowohl die Spec als auch den Code erzeugt, korrelieren die Fehler
(\emph{common-mode}) --- genau der von VeriAct \cite{veriact2026} beschriebene
Fall, dass verifizierer-akzeptiert nicht korrekt heißt. Der gefangene Mis-Spec
in \texttt{ledger-casestudy}\,@\,\texttt{0f39ab0} (\glqq Falsifiable, after 4 tests\grqq;
reproduzierbares Negativ-Artefakt \texttt{demo-gate-at-failure.sh}) wurde nicht
durch den Loop aus sich selbst entdeckt, sondern weil ein Mensch eine stärkere
Property und ein geändertes Requirement als Gegen-Orakel einbrachte. Der Mensch
ist hier der Dekorrelator; die Trennung Generator~$\perp$~Spec-Autor gilt nur,
wenn die Spec aus separater Quelle stammt.

\subsection{Was in der Spalte B bleibt}
\label{sec:grenzen-restb}

In Anlehnung an die A/B-Verifizierbarkeitstabelle (\Cref{sec:grenzen-ab})
markieren wir die folgenden Behauptungen ehrlich als~\AB{B}, also als menschliche
Setzung ohne maschinelle Außenprüfung:

\begin{itemize}
  \item die Bindung der Spec an die Welt (Spec~$\models$~Welt) --- nach
        VeriAct \cite{veriact2026} nicht orakel-prüfbar;
  \item das Setzen der Grund-Invariante (normativ, nicht berechnungstheoretisch;
        \Cref{sec:reflexiv-grenze});
  \item Neuheit und Abgrenzung --- eine Literatur-Aussage, kein Beweis
        (\Cref{sec:related});
  \item die Skalierung der Methode über Domänen --- $N=2$ zeigt Anwendbarkeit,
        nicht Skalierung; der lokale Lean-Allquantor über Buchungsfolgen
        ist~\AB{A}, der Sprung auf \glqq für jedes System\grqq{} ist~\AB{B}.
\end{itemize}

\subsection{A/B-Verifizierbarkeitstabelle}
\label{sec:grenzen-ab}
Jede tragende Behauptung dieser Arbeit deklariert sich als~\AB{A} (maschinell und
extern verifiziert: Commit-Hash, arXiv-ID, CI-Run, \texttt{sorry}-Audit) oder
als~\AB{B} (menschliche Setzung). \Cref{tab:ab} listet die tragenden Behauptungen mit
ihrer Marke und ihrem Anker auf. Keine tragende Assertion bleibt unmarkiert; das
ist der zentrale Hebel von Logik gegenüber bloßer LLM-Plausibilität.

\begin{table}[htbp]
  \centering
  \small
  \caption{A/B-Verifizierbarkeit jeder tragenden Behauptung. A~$=$~maschinell und
  extern verifiziert (Commit-Hash, arXiv-ID, CI-Run, \texttt{sorry}-Audit);
  B~$=$~menschliche Setzung.}
  \label{tab:ab}
  \begin{tabular}{@{}p{0.50\textwidth} c p{0.32\textwidth}@{}}
    \toprule
    \textbf{Behauptung} & \textbf{A/B} & \textbf{Anker} \\
    \midrule
    Runenruf: $46/46$ Tests grün & \AB{A} & \texttt{runenruf}\,@\,\texttt{56376a5} (Runner) \\
    Ledger: $6/6$ Tests grün & \AB{A} & \texttt{ledger-casestudy}\,@\,\texttt{0f39ab0} (Runner) \\
    CDD-Selbst-Gate: $37/37$ Tests, $62/66$ \texttt{Aligned} & \AB{A} & \texttt{cong-driven-development}\,@\,\texttt{main} (Runner) \\
    Werterhaltung + Nichtnegativität \texttt{sorry}-frei & \AB{A} & \texttt{\#print axioms} (nur \texttt{propext}, \texttt{Quot.sound})\,@\,\texttt{0f39ab0} \\
    Lean ohne Mathlib, CI \texttt{lean-proof} grün & \AB{A} & CI-Job \texttt{lean-proof}\,@\,\texttt{0f39ab0} \\
    Allowlist-Trennung Generator~$\perp$~Prüf-Orakel & \AB{A} & Code\,@\,\texttt{main} \\
    Negativ-Artefakt \texttt{demo-gate-at-failure.sh} & \AB{A} & Skript\,@\,\texttt{0f39ab0} \\
    FsCheck \texttt{spec-siegel-lager-nichtnegativ} (jeder Seed/Folge) & \AB{A} & \texttt{runenruf}\,@\,\texttt{56376a5} (Runner) \\
    Kompressionsfaktoren ($\sim\!40$--$125\times$) als LoC-Zahlen & \AB{A} & Code\,@\,\texttt{56376a5}/\texttt{0f39ab0} \\
    Zitate (vollständige Bib-Daten) & \AB{A} & arXiv-ID / Journalband \\
    \midrule
    Spec trifft die Welt (Spec~$\models$~Welt) & \AB{B} & per VeriAct nicht orakel-prüfbar \\
    \emph{Wer} die Grund-Invariante setzt (Mensch) & \AB{B} & normative Wahl, \Cref{sec:reflexiv-grenze} \\
    Stärkere Ledger-Property als Gegen-Orakel & \AB{B} & menschliche Setzung, \Cref{sec:ledger-misspec} \\
    Neuheit/Abgrenzung der Kombination & \AB{B} & Literatur-Aussage, \Cref{sec:related-neuheit} \\
    Methoden-Skalierung über Domänen & \AB{B} & $N=2$; Sprung auf \glqq jedes System\grqq{} \\
    \glqq Fixierte Eigenschaft ist die richtige\grqq{} & \AB{B} & Spec-gegen-Welt, \Cref{bem:kompression} \\
    \bottomrule
  \end{tabular}
\end{table}

\section{Verwandte Arbeiten}
\label{sec:related}

\subsection{Intent als offene Herausforderung}
\label{sec:related-intent}

Lahiri \cite{lahiri2026} formuliert die Formalisierung von Intent als benannte,
\emph{offene} Grand Challenge zuverlässigen Programmierens im Zeitalter von
KI-Agenten. Wir beanspruchen keine Lösung dieser Challenge, sondern eine
partielle, getypte Externalisierung des Intents pro Spec --- und markieren die
verbleibende Spec-gegen-Welt-Bindung als~\AB{B} (\Cref{sec:grenzen-restb}). Storey
\cite{storey2026} prägt den komplementären Begriff der \emph{Intent Debt},
\glqq the absence of externalized rationale that developers and AI agents need to
work safely with code\grqq; der externalisierte \texttt{.spot/}-Graph ist genau
ein Versuch, diese Schuld nicht aufzubauen.

\subsection{LLM im Modulo-Rahmen}
\label{sec:related-modulo}

Kambhampati und andere \cite{kambhampati2024} argumentieren, dass LLMs nicht
selbst planen, wohl aber im \emph{LLM-Modulo}-Rahmen mit externen Verifikatoren
nützlich sind. Unsere Architektur ist eine Instanz dieses Musters: Der Generator
schlägt vor, ein externer, mechanisch getrennter Prüfer entscheidet
(\Cref{sec:grenzen-commonmode}).

\subsection{Nächste Nachbarschaft}
\label{sec:related-veriact}

VeriAct \cite{veriact2026} ist die nächste Nachbarschaft: Die Arbeit zeigt, dass
verifizierer-akzeptierte Spezifikationen nicht korrekt sein müssen, und benennt
damit präzise die Grenze, an der unser menschlicher Dekorrelator ansetzt. He und
andere \cite{pgs2025} verfolgen mit \emph{Property-Oriented and Structurally
Minimal Feedback} ein verwandtes Ziel auf der Feedback-Seite: minimale,
eigenschaftsorientierte Rückmeldung zur Code-Verfeinerung. Beide Linien teilen
unsere Prämisse, dass das Orakel die kritische Stelle ist, nicht der Loop.

\subsection{Trajektorie und Werkzeuge}
\label{sec:related-trajektorie}

METR \cite{metr2025} misst den \glqq50\%-task-completion time horizon\grqq: Zu
Beginn 2025 liegt er bei etwa $50$ bis $110$~Minuten, mit einer Verdopplung etwa
alle $7$~Monate. Wir referenzieren ausschließlich diese Zahlen; sie motivieren,
warum das Terminierungs-Orakel mit wachsender Aufgabenlänge wichtiger wird, nicht
weniger. Werkzeuge wie Kiro und Spec-Kit verfolgen eine
spezifikationsgetriebene Arbeitsweise; sie verankern die Spec jedoch nicht in
einem typisierten Graphen mit architektonisch getrenntem Generator/Orakel und
property- bzw.\ beweisverifiziertem Gate.

\subsection{Das breitere Feld}
\label{sec:related-feld}

Spezifikationsgetriebene Entwicklung ist 2026 kein Randphänomen, sondern die
dominante Antwort auf \glqq Vibe-Coding\grqq: neben Kiro und Spec-Kit etwa Cursor,
OpenSpec, BMAD, Google Antigravity sowie die kommerzielle Plattform
Tessl~\cite{tessl2026}, die die Spec als primäres Artefakt führt und mit Tests als
Leitplanken paart --- ihr Problem-Framing (Agenten, die \glqq fertig\grqq{} melden,
ohne es zu sein) deckt sich mit unserer Diagnose. Parallel hat sich ein eigenes
Subfeld, \emph{vericoding}, etabliert: aus formalen Spezifikationen verifizierten
Code synthetisieren und mit Lean, Dafny oder Verus prüfen --- etwa
CLEVER~\cite{clever2025}, ein Benchmark für end-to-end in Lean verifizierte
Codegenerierung, oder VeriGuard~\cite{veriguard2025}, das Agenten-Sicherheit über
verifizierte Codegenerierung mit getrenntem Generator und Prüfer absichert. Unser
Beitrag liegt damit weder im Gedanken \glqq Spec als Wahrheit\grqq{} (heute
Commodity) noch in \glqq Code gegen formale Spec verifizieren\grqq{} (vericoding),
sondern in deren \emph{Integration} zu einer laufenden, getypten
Entwicklungsmethode: dem Konvergenz-Zustandsgraphen (\texttt{SPOT}), der
architektonischen Generator/Orakel-Trennung per Allowlist und der reflexiven
Selbstanwendung. Das ist eine Architektur- und Integrationsaussage, kein
Ideen-Primat.

\subsection{Neuheit, konditional}
\label{sec:related-neuheit}

Wir erheben keinen Prioritätsanspruch. Uns ist keine publizierte Implementierung
mit \emph{genau dieser Kombination} bekannt: getypte, externalisierte Spec;
architektonisch getrennter Generator und Prüf-Orakel; Werkzeug-Allowlist ohne
Schreibrecht des Generators; sowie ein property- und beweisverifiziertes Gate.
Diese Aussage ist eine Literatur-Aussage und damit~\AB{B}
(\Cref{sec:grenzen-restb}) --- kein Beweis.

\section{Fazit}
\label{sec:fazit}

Das eigentliche Problem agentischer Entwicklung ist nicht der Loop --- den zu
schreiben ist, wie Cherny \cite{cherny2026} es formuliert, die Routinearbeit ---
sondern sein Terminierungs-Orakel. Unsere Antwort ist Konvergenz gegen eine
externalisierte, getypte Spec, geprüft durch ein vom Generator architektonisch
getrenntes Orakel. Wir belegen das mit zwei kleinen Fallstudien
(\texttt{runenruf}\,@\,\texttt{56376a5}, $46/46$; \texttt{ledger-casestudy}\,@\,\texttt{0f39ab0}, $6/6$),
einem reflexiven Selbst-Gate (\texttt{cong-driven-development}\,@\,\texttt{main},
$37/37$, $62/66$ \texttt{Aligned}) und genau zwei maschinell bewiesenen
Invarianten (\texttt{Werterhaltung.lean}: \texttt{werterhaltung} + \texttt{nichtnegativ}, \texttt{sorry}-frei).

Wir schließen mit einer Lücke, nicht mit einem Versprechen. Die Spec-gegen-Welt-Bindung
bleibt~\AB{B}; die Grund-Invariante terminiert den Regress nur durch einen externen
Anker, dessen Setzung eine Wahl ist und keine berechnungstheoretische Notwendigkeit;
die Beweiskette ist erst für eine einzige Eigenschaft geschlossen; $N=2$ zeigt
Anwendbarkeit, nicht Skalierung. Aus keinem Theorem folgt, dass ein Mensch hier
unersetzlich wäre --- ein zweiter, hinreichend unabhängiger Agent könnte den
Dekorrelator stellen. Was wir zeigen, ist bescheidener und überprüfbarer: dass
das Orakel die richtige Stelle ist, um anzusetzen, und dass eine Schraube weiter
gedreht ist --- mit Lücke, ehrlich markiert.

\appendix

\section{Reproduzierbarkeit}
\label{app:repro}

Alle in \Cref{tab:ab} als~\AB{A} markierten Behauptungen sind an genau drei Repositories
zu festen Commits gebunden. \Cref{tab:repos} fasst die Anker zusammen; die
nachfolgenden Befehle reproduzieren die autoritativen Runner-Counts, den
\texttt{sorry}-freien Lean-Lauf und das Negativ-Artefakt.

\begin{table}[htbp]
  \centering
  \small
  \caption{Verifizierte Artefakte und ihre Anker (Repository\,@\,Commit).}
  \label{tab:repos}
  \begin{tabular}{@{}lll@{}}
    \toprule
    \textbf{Repository} & \textbf{Commit} & \textbf{Rolle} \\
    \midrule
    \texttt{Koschnag/cong-driven-development} & \texttt{main} & Methode/IDE, reflexives Selbst-Gate \\
    \texttt{Koschnag/runenruf} & \texttt{56376a5} & Existenzbeweis (Simulationsdomäne) \\
    \texttt{Koschnag/ledger-casestudy} & \texttt{0f39ab0} & Gegenstück (echtes IO, Lean-Beweis) \\
    \bottomrule
  \end{tabular}
\end{table}

\subsection{Methode und Selbst-Gate}
\label{app:repro-cdd}

\begin{lstlisting}[language=bash,caption={CDD-Selbst-Gate: $37/37$ Tests, $62/66$ \texttt{Aligned}.},label={lst:repro-cdd}]
git clone https://github.com/Koschnag/cong-driven-development.git
cd cong-driven-development
git checkout main
dotnet test tests/Cdd.Tests   # autoritativer Runner-Count: 37/37
\end{lstlisting}

\subsection{Fallstudie A: Runenruf}
\label{app:repro-runenruf}

\begin{lstlisting}[language=bash,caption={Runenruf: $46/46$ Tests, FsCheck \texttt{spec-siegel-lager-nichtnegativ}.},label={lst:repro-runenruf}]
git clone https://github.com/Koschnag/runenruf.git
cd runenruf
git checkout 56376a5
dotnet test tests/Runenruf.Tests   # autoritativer Runner-Count: 46/46
\end{lstlisting}

\subsection{Fallstudie B: Ledger (Tests, Lean, Negativ-Artefakt)}
\label{app:repro-ledger}

\begin{lstlisting}[language=bash,caption={Ledger: $6/6$ Tests, \texttt{sorry}-freier Lean-Beweis, reproduzierbares Negativ-Artefakt.},label={lst:repro-ledger}]
git clone https://github.com/Koschnag/ledger-casestudy.git
cd ledger-casestudy
git checkout 0f39ab0

# 1) F#-Tests (autoritativer Runner-Count: 6/6)
dotnet test tests/Ledger.Tests

# 2) Lean-Beweis (sorry-frei; #print axioms: nur propext, Quot.sound)
lean proofs/Werterhaltung.lean

# 3) reproduzierbares NEGATIV-Artefakt: Gate im Falsifikationszustand
#    ("Falsifiable, after 4 tests")
./demo-gate-at-failure.sh
\end{lstlisting}

\noindent
Der Befehl \texttt{lean proofs/Werterhaltung.lean} aus \Cref{lst:repro-ledger}
schließt den Beweis ohne Mathlib (nur \texttt{Core}~$+$~\texttt{omega}) ab;
\texttt{\#print axioms werterhaltung} weist ausschließlich \texttt{propext} und
\texttt{Quot.sound} aus. Das Skript \texttt{demo-gate-at-failure.sh} ist kein
Positiv-, sondern ein \emph{Negativ}-Artefakt: Es führt das Gate reproduzierbar in
den Zustand \glqq Falsifiable, after 4 tests\grqq{} und macht damit den in
\Cref{sec:ledger-misspec} beschriebenen, vom Menschen gefangenen Mis-Spec
wiederholbar sichtbar.

\begin{thebibliography}{10}

\bibitem{lahiri2026}
Lahiri, Shuvendu K.:
\newblock Intent Formalization: A Grand Challenge for Reliable Coding in the Age
of AI Agents.
\newblock \texttt{arXiv:2603.17150}, 2026.

\bibitem{storey2026}
Storey, Margaret-Anne:
\newblock From Technical Debt to Cognitive and Intent Debt: Rethinking Software
Health in the Age of AI.
\newblock \texttt{arXiv:2603.22106}, 2026.

\bibitem{kambhampati2024}
Kambhampati, Subbarao and others:
\newblock LLMs Can't Plan, But Can Help Planning in LLM-Modulo Frameworks.
\newblock \texttt{arXiv:2402.01817}, 2024.

\bibitem{veriact2026}
Misu, Md Rakib Hossain; Ma, Iris; Lopes, Cristina V.:
\newblock VeriAct: Beyond Verifiability --- Agentic Synthesis of Correct and
Complete Formal Specifications.
\newblock \texttt{arXiv:2604.00280}, 2026.

\bibitem{pgs2025}
He, Lehan; Chen, Zeren; Zhang, Zhe; Gao, Xiang; Sheng, Lu:
\newblock Effective LLM Code Refinement via Property-Oriented and Structurally
Minimal Feedback.
\newblock \texttt{arXiv:2506.18315}, 2025.

\bibitem{metr2025}
METR:
\newblock Measuring AI Ability to Complete Long Software Tasks.
\newblock \texttt{arXiv:2503.14499}, 2025.

\bibitem{rice1953}
Rice, Henry Gordon:
\newblock Classes of Recursively Enumerable Sets and Their Decision Problems.
\newblock Transactions of the American Mathematical Society 74(2):358--366, 1953.

\bibitem{cherny2026}
Cherny, Boris:
\newblock zitiert in \glqq Loop Engineering\grqq, The New Stack, 2026 ---
\glqq my job is to write loops\grqq.

\bibitem{cdd2026}
Nguyen, Cong Chanh Vinzenz:
\newblock cong-driven-development: A Spec-Convergent Method and IDE for Agentic
Software Development.
\newblock Software-Repository, \texttt{github.com/Koschnag/cong-driven-development}\,@\,\texttt{main}, 2026.

\bibitem{tessl2026}
Tessl:
\newblock Spec-Driven Development with Tessl (Framework und Spec-Registry).
\newblock Agent-Enablement-Plattform, \url{https://tessl.io}, 2026.

\bibitem{clever2025}
Thakur, Amitayush u.\,a.:
\newblock CLEVER: A Curated Benchmark for Formally Verified Code Generation.
\newblock arXiv:2505.13938, 2025.

\bibitem{veriguard2025}
Miculicich, Lesly u.\,a.:
\newblock VeriGuard: Enhancing LLM Agent Safety via Verified Code Generation.
\newblock arXiv:2510.05156, 2025.

\end{thebibliography}

\end{document}

